% -*- coding: utf-8-unix; time-stamp-format: "%02d-%02m-%:y %02H:%02M:%02S %Z"; sentence-end-double-space: t; -*-
%<*dtx>
\def\bnedtxtimestamp  {Time-stamp: <17-05-2021 14:16:23 CEST>}
\iffalse
%</dtx>
%<*drv>
%% ---------------------------------------------------------------
\def\bnedocdate {2021/05/17}% package bnumexpr documentation date
\def\bnepackdate{2021/05/17}% package bnumexpr date
\def\bneversion {1.5}      % package bnumexpr version
%</drv>
%<*readme>--------------------------------------------------------
| Source:  bnumexpr.dtx
| Version: v1.5, 2021/05/17 (doc: 2021/05/17)
| Author:  Jean-Francois Burnol
| Info:    Expressions with big integers
| License: LPPL 1.3c

bnumexpr usage
==============

The LaTeX package `bnumexpr` allows expandable computations with
integers and the four infix operators `+`, `-`, `*`, `/` using the
expression syntax familiar from the `\numexpr` e-TeX parser, with
these extensions:

- arbitrarily big integers,
- floored division `//`,
- associated modulo `/:`,
- power operators `^` and `**`,
- factorial post-fix operator `!`,
- comma separated expressions,
- the space character as well as the underscore may serve
  to separate groups of digits,
- optional conversion of output to hexadecimal,
- customizability and extendibility of the syntax.

The expression parser is a scaled-down variant from the
`\xintiiexpr...\relax`
parser from package [xintexpr](http://ctan.org/pkg/xintexpr).

To support hexadecimal input and output, the package
[xintbinhex](http://ctan.org/pkg/xint) is loaded automatically.

The package loads by default [xintcore](http://ctan.org/pkg/xint)
but the option _custom_ together with macro `\bnumexprsetup` allow to map
the syntax elements to macros from an alternative big integer
expandable engine of the user own choosing,
and then [xintcore](http://ctan.org/pkg/xint) is not loaded.

Installation
============

Use your installation manager to install or update `bnumexpr`.

Else, obtain `bnumexpr.dtx`, from CTAN:

> <http://www.ctan.org/pkg/bnumexpr>

Run `"etex bnumexpr.dtx"` to extract these files:

`bnumexpr.sty`
  : this is the style file.

`README.md`

`bnumexprchanges.tex`
  : change history.

`bnumexpr.tex`
  : can be used to generate the documentation

To generate the documentation:

  - with latex+dvipdfmx: `"latex bnumexpr.tex"` (thrice) then
    `"dvipdfmx bnumexpr.dvi"`.

  - with pdflatex: `"pdflatex bnumexpr.tex"` (thrice).

In both cases files `README.md` and `bnumexprchanges.tex` must
be located in the same repertory as `bnumexpr.tex` and `bnumexpr.dtx`.

Without `bnumexpr.tex`:

  - `"pdflatex bnumexpr.dtx"` (thrice) extracts all files and
    simultaneously generates the pdf documentation.

Final steps:

- move files to appropriate destination:

           bnumexpr.sty   --> TDS:tex/latex/bnumexpr/

           bnumexpr.dtx   --> TDS:source/latex/bnumexpr/

           bnumexpr.pdf   --> TDS:doc/latex/bnumexpr/
              README.me   --> TDS:doc/latex/bnumexpr/

- discard auxiliary files generated during compilation.

License
=======

Copyright (C) 2014-2021 by Jean-Francois Burnol

| This Work may be distributed and/or modified under the
| conditions of the LaTeX Project Public License 1.3c.
| This version of this license is in

>   <http://www.latex-project.org/lppl/lppl-1-3c.txt>

| and version 1.3 or later is part of all distributions of
| LaTeX version 2005/12/01 or later.

This Work has the LPPL maintenance status "author-maintained".

The Author and Maintainer of this Work is Jean-Francois Burnol.

This Work consists of the main source file and its derived files

    bnumexpr.dtx, bnumexpr.sty, bnumexpr.pdf, bnumexpr.tex,
    bnumexprchanges.tex, README.md

%</readme>--------------------------------------------------------
%<*!readme>
%% ---------------------------------------------------------------
%% The bnumexpr package: Expressions with big integers
%% Copyright (C) 2014-2021 by Jean-Francois Burnol
%%
%</!readme>
%<*changes>-------------------------------------------------------
\item[1.5 (2021/05/17)]
  \begin{itemize}
  \item \textbf{breaking change:} the power operators act now in a right
    associative way; this has been announced at \xintexprname as a probable
    future evolution, and is implemented in anticipation here now.
  \item \textbf{fix two bugs} (imported from upstream \xintexprname) regarding
    hexadecimal input: impossibility to use |"\foo| syntax (one had to do
    |\expandafter"\foo| which is unexpected constraint; a very longstanding
    \xintexprname bug) and issues with leading zeros (since \xintexprname |1.2m|).
  \item renamed |\bnumexprsetup| into |\bnumsetup|; the former remains
    available but is deprecated.
  \item the customizability and extendibility is now total:
    \begin{enumerate}
      \item |\bnumprintone|, |\bnumprintonetohex|, |\bnumprintonesep|, |\bnumhextodec|,
      \item |\bnumdefinfix| which allows to add extra infix operators,
      \item |\bnumdefpostfix| which allows to add extra postfix operators.
    \end{enumerate}
  \item |\bnumsetup|, |\bnumdefinfix|, |\bnumdefpostfix| obey the
    |\xintglobaldefstrue| and |\xintverbosetrue| settings.
  \item documentation is extended, providing details regarding the precedence
    model of the parser, as inherited from upstream \xintexprname; also an
    example of usage of |\bnumsetup| is included on how to transform
    |\bnumeval| into a calculator with fractions.
  \end{itemize}

\item[1.4a (2021/05/13)]
  \begin{itemize}
  \item fix undefined control sequences errors encountered by the parser in
    case of either extra or missing closing parenthesis (due to a problem
    in technology transfer at |1.4| from upstream \xintexprname).
  \item fix |\BNE_Op_opp| must now be \emph{f}-expandable (also caused as a
    collateral to the technology transfer).
  \item fix user documentation regarding the constraints applying to the user
    replacement macros for the core algebra, as they have changed at |1.4|.
  \end{itemize}

\item[1.4 (2021/05/12)]
  \begin{itemize}
  \item technology transfer from \xintexprname |1.4| of
    |2020/01/31|. The |\expanded| primitive is now required (TeXLive
    2019).
  \item addition to the syntax of the |"| prefix for hexadecimal
    input.
  \item addition of |\evaltohex| which is like |\bnumeval| with an extra
    conversion step to hexadecimal notation.
  \end{itemize}

\item[1.2e (2019/01/08)]
  Fixes a documentation glitch (extra braces when mentioning |\the\numexpr| or
  |\thebnumexpr|).

\item[1.2d (2019/01/07)]
  \begin{itemize}
  \item requires \xintcorename |1.3d| or later (if not using option |custom|).
  \item adds |\bnumeval|\marg{expression} user interface.
  \end{itemize}

\item[1.2c (2017/12/05)] \textbf{Breaking changes:}
  \begin{itemize}
  \item requires \xintcorename |1.2p| or later (if not using option |custom|).
  \item |divtrunc| key of |\bnumexprsetup| is renamed to |div|.
  \item the |//| and |/:| operators are now by default associated to the
    \emph{floored} division. This is to keep in sync with the change of
    \xintcorename at |1.2p|.
  \item for backwards compatibility, one
    may add to existing document:\newline
    \string\bnumexprsetup\{div=\string\xintiiDivTrunc,
    mod=\string\xintiiModTrunc\}
  \end{itemize}

\item[1.2b (2017/07/09)]
  \begin{itemize}
  \item the |_| may be used to separate visually blocks of digits in long
    numbers.
  \end{itemize}

\item[1.2a (2015/10/14)]
  \begin{itemize}
  \item requires \xintcorename |1.2| or later (if not using option |custom|).
  \item additions to the syntax: factorial |!|, truncated division
    |//|, its associated modulo |/:| and |**| as alternative to |^|.
  \item all options removed except |custom|.
  \item new command |\bnumexprsetup| which replaces the commands such as
    |\bnumexprusesbigintcalc|.
  \item the parser is no more limited to numbers with at most 5000 digits.
  \end{itemize}

\item[1.1b (2014/10/28)]
  \begin{itemize}
  \item README converted to |markdown/pandoc| syntax,
  \item the package now loads only |xintcore|, which belongs to |xint| bundle
    version |1.1| and extracts from the earlier |xint| package the core arithmetic
    operations as used by |bnumexpr|.
  \end{itemize}

\item[1.1a (2014/09/22)]
  \begin{itemize}
  \item added |l3bigint| option to use experimental \LaTeX3
    package of the same name,

  \item added Changes and Readme sections to the documentation,

  \item better |\BNE_protect| mechanism for use of
  |\bnumexpr...\relax| inside an |\edef| (without |\bnethe|). Previous one,
  inherited from |xintexpr.sty 1.09n|, assumed that the |\.=<digits>| dummy
  control sequence encapsulating the computation result had |\relax|
  meaning. But removing this assumption was only a matter of letting
  |\BNE_protect| protect two, not one, tokens. This will be backported to
  next version of \xintexprname, naturally (done with |xintexpr.sty 1.1|).
  \end{itemize}

\item[1.1 (2014/09/21)] First release. This is down-scaled from the
  (development version of) \xintexprname. Motivation came the previous day
  from a chat with \textsc{Joseph Wright} over big int status in \LaTeX3.
  The |\bnumexpr...\relax| parser can be used on top of big int macros of
  one's choice. Functionalities limited to the basic operations. I leave
  the power operator |^| as an option.
%</changes>-------------------------------------------------------
%<*drv>-----------------------------------------------------------
%% latex bnumexpr.tex (thrice) && dvipdfmx bnumexpr.dvi
%% to produce bnumexpr.pdf
%%
%% or pdflatex bnumexpr.tex (no need to change \Withdvipdfmx toggle,
%% pdf engine will be detected automatically)
%%
\NeedsTeXFormat{LaTeX2e}
\ProvidesFile{bnumexpr.tex}%
[\bnepackdate\space v\bneversion\space driver file for %
  bnumexpr documentation (JFB)]%
\PassOptionsToClass{a4paper,fontsize=11pt}{scrdoc}
\chardef\Withdvipdfmx 1 %
\chardef\NoSourceCode 0 % replace 0 by 1 for not including source code
\input bnumexpr.dtx
%%% Local Variables:
%%% mode: latex
%%% End:
%</drv>-----------------------------------------------------------
%<*dtx>
\fi % end of \iffalse block around generated files
\chardef\noetex 0
\ifx\numexpr\undefined \chardef\noetex 1 \fi
\ifnum\noetex=1 \chardef\extractfiles 0 % extract files, then stop
\else
    \ifx\ProvidesFile\undefined
      \chardef\extractfiles 0 % no LaTeX2e; etex, ... on bnumexpr.dtx
    \else % latex/pdflatex on bnumexpr.tex or on bnumexpr.dtx
      \ifx\Withdvipdfmx\undefined
        % latex run is on bnumexpr.dtx, we will extract all files
        \chardef\extractfiles 1 % 1 = extract and typeset, 2=only typeset
        \chardef\Withdvipdfmx 0 % 0 = pdflatex or latex+dvips, 1 = dvipdfmx
        \chardef\NoSourceCode 0 % 0 = include source code, 1 = do not
        \NeedsTeXFormat{LaTeX2e}%
        \PassOptionsToClass{a4paper,fontsize=11pt}{scrdoc}%
      \else % latex run is on bnumexpr.tex,
        \chardef\extractfiles 2 % no extractions
      \fi
      \ProvidesFile{bnumexpr.dtx}[bundle source (\bnedtxtimestamp)]%
    \fi
\fi
\ifnum\extractfiles<2 % extract files
\def\MessageDeFin{\newlinechar10 \let\Msg\message
\Msg{^^J}%
\Msg{********************************************************************^^J}%
\Msg{*^^J}%
\Msg{* To finish the installation you have to move the following^^J}%
\Msg{* files into a directory searched by TeX:^^J}%
\Msg{*^^J}%
\Msg{*\space\space\space\space bnumexpr.sty^^J}%
\Msg{*^^J}%
\Msg{* To produce the documentation run latex thrice on bnumexpr.tex^^J}%
\Msg{* then dvipdfmx on bnumexpr.dvi. (ignore the dvipdfmx warnings)^^J}%
\Msg{*^^J}%
\Msg{* Happy TeXing!^^J}%
\Msg{*^^J}%
\Msg{********************************************************************^^J}%
}%
\begingroup
    \input docstrip.tex
    \askforoverwritefalse
    \generate{\nopreamble\nopostamble
    \file{README.md}{\from{bnumexpr.dtx}{readme}}
    \usepreamble\defaultpreamble
    \usepostamble\defaultpostamble
    \file{bnumexprchanges.tex}{\from{bnumexpr.dtx}{changes}}
    \file{bnumexpr.tex}{\from{bnumexpr.dtx}{drv}}
    \file{bnumexpr.sty}{\from{bnumexpr.dtx}{package}}}
\endgroup
\fi % end of file extraction
\ifnum\extractfiles=0
% direct tex/etex/xetex on bnumexpr.dtx, files now extracted, stop
  \MessageDeFin\expandafter\end
\fi
% no use of docstrip to extract files if latex compilation was on bnumexpr.tex
\ifdefined\MessageDeFin\AtEndDocument{\MessageDeFin}\fi
%-------------------------------------------------------------------------------
\documentclass {scrdoc}
\ifnum\NoSourceCode=1 \OnlyDescription\fi
\usepackage{ifpdf}
\ifpdf\chardef\Withdvipdfmx 0 \fi
\makeatletter
\ifnum\Withdvipdfmx=1
   \@for\@tempa:=hyperref,bookmark,graphicx,xcolor,pict2e\do
            {\PassOptionsToPackage{dvipdfmx}\@tempa}
   %
   \PassOptionsToPackage{dvipdfm}{geometry}
   \PassOptionsToPackage{bookmarks=true}{hyperref}
   \PassOptionsToPackage{dvipdfmx-outline-open}{hyperref}
   %\PassOptionsToPackage{dvipdfmx-outline-open}{bookmark}
   %
   \def\pgfsysdriver{pgfsys-dvipdfm.def}
\else
   \PassOptionsToPackage{bookmarks=true}{hyperref}
\fi
\makeatother

\pagestyle{headings}

\usepackage[T1]{fontenc}

\usepackage[hscale=0.66,vscale=0.75]{geometry}

% requires newtxtt 1.05 or later
\usepackage[zerostyle=a,scaled=0.95,straightquotes]{newtxtt}
\renewcommand\familydefault\ttdefault
\usepackage[noendash]{mathastext}
\renewcommand\familydefault\sfdefault

\usepackage{graphicx}
\usepackage[dvipsnames]{xcolor}
\definecolor{joli}{RGB}{225,95,0}
\definecolor{JOLI}{RGB}{225,95,0}
\definecolor{BLUE}{RGB}{0,0,255}
%\definecolor{niceone}{RGB}{38,128,192}% utilisé avant pour urlcolor
%\colorlet{smallverbcolor}{RoyalPurple}% \colorlet{jfverbcolor}{yellow!5}
\colorlet{smallverbcolor}{NavyBlue}%
\colorlet{softwrapcolor}{blue}
\colorlet{digitscolor}{OrangeRed}

% transféré de xint-manual.tex (maintenant dans xint.dtx)
\DeclareFontFamily{U}{MdSymbolC}{}
\DeclareFontShape {U}{MdSymbolC}{m}{n}{<-> MdSymbolC-Regular}{}

\makeatletter
\newbox\cdbx@SoftWrapIcon
\def\cdbx@SetSoftWrapBox{%
    \setbox\cdbx@SoftWrapIcon\hb@xt@\z@
    {\hb@xt@\fontdimen2\font
        {\hss{\color{softwrapcolor}\usefont{U}{MdSymbolC}{m}{n}\char"97}\hss}%
     \hss}%
   }
\makeatother

\usepackage[english]{babel}

\usepackage{hyperref}
\hypersetup{%
%linktoc=all,%
breaklinks=true,%
colorlinks=true,%
urlcolor=SkyBlue,%
linkcolor=PineGreen,%
pdfauthor={Jean-François Burnol},%
pdftitle={The bnumexpr package},%
pdfsubject={Arithmetic with TeX},%
pdfkeywords={Expansion, arithmetic, TeX},%
pdfstartview=FitH,%
pdfpagemode=UseOutlines}
\usepackage{bookmark}

% importé de xint 2021/05/14
\newcommand\RaisedLabel[2][6]{%
\vspace*{-#1\baselineskip}%
\begingroup
  \let\leavevmode\relax\phantomsection
  \label{#2}%
\endgroup
\vspace*{#1\baselineskip}%
}

%---- \verb, and verbatim like `environments'. \MicroFont et \MacroFont
\makeatletter
\def\MicroFont {\ttfamily\cdbx@SetSoftWrapBox\color{smallverbcolor}}
%\def\MicroFont {\ttfamily \color[named]{OrangeRed}\cdbx@SetSoftWrapBox }

% \MacroFont est utilisé par macrocode, mais sa définition est figée dans
% \macro@font au \begin{document}

\def\MacroFont {\ttfamily \baselineskip12pt\relax }

%--- November 4, 2014 making quotes straight. (maintenu pour *)
%
% there is no hook in \macrocode after \dospecials etc. Thus I will need to
% take the risk that some future evolution of doc.sty (or perhaps scrdoc)
% invalidates the following.
%
% Actually, I should not at all rely on the doc class, I should do it all by
% myself.

\def\macrocode{\macro@code
               \frenchspacing \@vobeyspaces
               \makestarlowast % \makequotesstraight (obsolète)
               \xmacro@code }

%--- lower * symbol in text

\def\lowast{\raisebox{-.25\height}{*}}
\catcode`* \active
\def\makestarlowast {\let*\lowast\catcode`\*\active}%
\catcode`* 12

\def\verb #1%
{%
  \relax\leavevmode\null
  \begingroup
   \MicroFont
   \def\@makeletter##1{\catcode`##1=11\relax}%
   % \scantokens will have a result of inserting a space after cs's.
   % hence the need to have the catcodes of things like _ right.
   % I also need < for > for code comments
   % \dospecials at begin document
   % \do\ \do\\\do\{\do\}\do\$\do\&\do\#\do\^\do\_\do\%\do\~\do\|
   \odef\dospecials{\dospecials\do\:\do\<\do\>\do\-\do\+}%
   % naturally won't work in footnotes though.
   % this code is truly not satisfying, but enough for my needs here.
   \catcode`\\ 11 \catcode`## 11
   % I don't do \catcode`\% 11 to avoid surprises in code comments
   % if my |...| has a linebreak
   \def\@jfverb ##1#1{\let\do\@makeletter \dospecials
   % lowering a bit the *
                      \makestarlowast
   %\let\do\do@noligs  \verbatim@nolig@list % not needed here
                      \@vobeyspaces\everyeof{\noexpand}%
                      \expandafter\@@jfverb\scantokens{##1}\relax}%
  \@jfverb
}%

% Pour v1.2 je laisse tomber les colorbox. À la place je change la couleur
% du texte. J'utilise le violet de xint. (mis dans \Microfont)
% \def\@@@jfverb #1{\ifcat\noexpand#1$% $
%                       \endgroup\else
% %      \penalty\z@
%        \discretionary{\copy\cdbx@SoftWrapIcon}{}{}%
% %      \colorbox{jfverbcolor}{\strut #1}% unfortunately, impacts pdf size
%        #1%
%        \expandafter\@@@jfverb\fi }

\def\@@jfverb #1{\ifcat\noexpand#1\relax\endgroup\else
       \discretionary{\copy\cdbx@SoftWrapIcon}{}{}%
       #1\expandafter\@@jfverb\fi }

\makeatother

% Lundi 17 mai 2021 à 13:58:27
% J'ajoute un export automatique des exemples qui servira de poor man test
% suite
\newwrite\tests
\immediate\openout\tests=\jobname-out.txt

% mardi 13 octobre 2015 à 15:32:51
% le | n'est pas encore actif ici!
\catcode`| \active
% first one uses \bnumexpr deliberately with no \bnethe to test 1.4
    \def\bneshow #1%
    {|\thebnumexpr#1\relax|%
     \edef\x{\bnumexpr#1\relax}%
     \immediate\write\tests{\detokenize{\thebnumexpr#1\relax}^^J\x}%
     $$\color{digitscolor}\x$$\ignorespaces}
    \def\bneshowthe #1%
    {|\bnethe\bnumexpr#1\relax|%
     \edef\x{\bnethe\bnumexpr#1\relax}%
     \immediate\write\tests{\detokenize{\bnethe\bnumexpr#1\relax}^^J\x}%
     $$\color{digitscolor}\x$$\ignorespaces}
    \def\bneshoweval #1%
    {|\bnumeval{#1}|%
     \edef\x{\bnumeval{#1}}%
     \immediate\write\tests{\detokenize{\bnumeval{#1}}^^J\x}%
     $$\color{digitscolor}\x$$\ignorespaces}
    \def\bneshowevaltohex #1%
    {|\evaltohex{#1}|%
     \edef\x{\evaltohex{#1}}%
     \immediate\write\tests{\detokenize{\evaltohex{#1}}^^J\x}%
     $$\color{digitscolor}\x$$\ignorespaces}
\catcode`| 12

% toujours pénible de devoir se battre avec les espaces verticaux de TeX
\AtBeginDocument{%
% \message{ICI below: \the\belowdisplayskip^^J%
%              belowshort: \the\belowdisplayshortskip^^J%
%              above: \the\abovedisplayskip^^J%
%              aboveshor: \the\abovedisplayshortskip^^J%
% }%
   \belowdisplayskip\belowdisplayshortskip
   \abovedisplayskip\abovedisplayshortskip
}

\DeclareRobustCommand\csa [1]{{\char92\detokenize{#1}}}
\newcommand\csh[1]{\texorpdfstring{\csa{#1}}{\textbackslash\detokenize{#1}}}

\usepackage{xspace}

\def\bnumname
   {\href{http://www.ctan.org/pkg/bnumexpr}{bnumexpr}\xspace }%

\def\bnumnameimp
   {\texorpdfstring
          {\hyperref[sec:bnumexprcode]{{\color{smallverbcolor}\ttzfamily bnumexpr}}}
          {bnumexpr}%
    \xspace }%

\def\bnumnameuserdoc
   {\texorpdfstring
          {\hyperref[titlepage]{{\color{PineGreen}\ttzfamily bnumexpr}}}
          {bnumexpr}%
    \xspace }%

\def\xintkernelname
   {\href{http://www.ctan.org/pkg/xint}{xintkernel}\xspace }%
\def\xintcorename
   {\href{http://www.ctan.org/pkg/xint}{xintcore}\xspace }%
\def\xintname
   {\href{http://www.ctan.org/pkg/xint}{xint}\xspace }%
\def\xintfracname
   {\href{http://www.ctan.org/pkg/xint}{xintfrac}\xspace }%
\def\xintexprname
   {\href{http://www.ctan.org/pkg/xintexpr}{xintexpr}\xspace }%
\def\xintbinhexname
   {\href{http://www.ctan.org/pkg/xint}{xintbinhex}\xspace }%

\frenchspacing
\usepackage{bnumexpr}
\usepackage{xintfrac}
\usepackage{etoc}
\usepackage{centeredline}
\usepackage{framed}
\makeatletter\let\check@percent\relax\makeatother
\begin{document}
\thispagestyle{empty}
\ttzfamily
\pdfbookmark[1]{Title page}{TOP}

{%
\normalfont\Large\parindent0pt \parfillskip 0pt\relax
 \leftskip 2cm plus 1fil \rightskip 2cm plus 1fil
 The \bnumname package\par
}

\RaisedLabel{titlepage}

{\centering
  \textsc{Jean-François Burnol}\par
  \footnotesize
  jfbu (at) free (dot) fr\par
  Package version: \bneversion\ (\bnepackdate);
            documentation date: \bnedocdate.\par
  {From source file \texttt{bnumexpr.dtx}. \bnedtxtimestamp.}\par
}

\etocsetnexttocdepth{section}
\tableofcontents

\section{\csh{bnumeval} (\csh{thebnumexpr}), \csh{evaltohex}}
\label{sec:bnumexpr}

\LaTeX\ Package \bnumname provides |\thebnumexpr|\meta{expression}|\relax|:
it is analogous to
  |\the\numexpr|\meta{expression}|\relax|, with these extensions:
\begin{itemize}
\item it allows arbitrarily big integers,
\item it computes powers (with either |**| or |^| as infix operator),
\item it computes factorials (with |!| as postfix operator),
\item it has an operator |//| for floored division and |/:| for the
  associated modulo (like {\catcode`\% 12 |%|} in Python which we can't use for obvious reasons),
\item the space character is ignored%
%
\footnote{It is not completely ignored, |\count37<space>| will automatically
  be prefixed by |\number| and the space token delimits the integer indexing
  the count register. Also, devious inputs using nested braces around spaces
  may create unexpected internal situations and even break the parser.}
%
and can thus be used to separate in the
  source blocks of digits for better readability of long numbers,
\item also the underscore |_| may be used as visual digit
  separator,
\item comma separated expressions are allowed,
\item syntax is customizable and extendible.
\end{itemize}
There is also a more core-level |\bnumexpr...\relax| construct%
\footnote{Since |1.4|, one can use |\bnumexpr...\relax| directly in
  typesetting context, it is not mandatory to prefix it with |\bnethe| or to
  use |\thebnumexpr|.}%
, which expands
to a self-contained unit, rather than to explicit digit tokens (and commas).
See \autoref{sec:differences} for some related information.

There is also the alternative interface |\bnumeval|\marg{expression}, where the
expression is fetched as braced argument.

And there is |\evaltohex|\marg{expression} which does the same as |\bnumeval|
but with a conversion to hexadecimal notation of the (possibly comma
separated) output. Hexadecimal input uses the |"| prefix.

This package parser is a scaled-down variant of |\xintiiexpr| from package
\xintexprname, dropping support for nested structures, functions, variables,
booleans, etc..., but incorporating by default support for hexadecimal input
as \xintbinhexname will be automatically loaded.

The \eTeX{} extensions are required, this is the default on all modern
installations for \verb?latex|pdflatex? and also for \verb?xelatex|lualatex?.

Further, at |1.4| (|2021/05/12|) the |\expanded| primitive is
required. It is available in all engines since \TeX Live 2019.

\section{Examples}

With certain languages, Babel with PDF\LaTeX\ may make some characters active, for example the
|!| with the French language. It must then be input as |\string!|.
\noindent\bneshow {---1 208 637 867 * (2 187 917 891 - 3 109 197 072)}%
\bneshowthe {(13_8089_1090-300_1890_2902)*(1083_1908_3901-109_8290_3890)}%
\bneshoweval {(92_874_927_979**5-31_9792_7979**6)/30!}%
\bneshoweval {30!/20!/21/22/23/24/25/(26*27*28*29)}%
\bneshoweval {13^50//12^50, 13^50/:12^50}%
\bneshoweval {13^50/12^50, 12^50}%
\bneshoweval {(1^10+2^10+3^10+4^10+5^10+6^10+7^10+8^10+9^10)^3}%
\bneshoweval {100!/36^100}%
\bneshoweval {"0010*"0100*"1000*"A0000, 16^(1+2+3+4)*10}%
\bneshowevaltohex{"7FFFFFFF+1, "0400^3, "ABCDEF*"0000FEDCBA, 1234}%
\bneshoweval {"\evaltohex{12345678}FFFF, 000012345679*16**4-1}

\section{The \texttt{custom} package option and \csh{bnumsetup}}

Package \bnumname needs that some \emph{big integer engine} provides the macros doing
the actual computations. By default, it loads package \xintcorename (a subset
of \href{http://www.ctan.org/pkg/xint}{xint}) and uses \csa{bnumsetup} in
the following way:
\begin{verbatim}
    \usepackage{xintcore}
    \bnumsetup{add=\xintiiAdd, sub=\xintiiSub, mul=\xintiiMul,
               divround=\xintiiDivRound, div=\xintiiDivFloor,
               mod=\xintiiMod, pow=\xintiiPow, fac=\xintiiFac,
               opp=\xintiiOpp}
\end{verbatim}

If
using \csa{bnumsetup}, it is not necessary to specify all keys,
for example one can do |\bnumsetup{mul=\MySlowerMul}|, and only
multiplication will be changed.

Naturally it is up to the user to load the appropriate package for
the alternative macros.

The macros serving as custom user replacements%
%
\footnote{%
  The replacement macros will by default receive arguments composed of
  explicit digit tokens, with no leading zeros, with at most one leading minus
  sign and no plus sign.

  The format of these arguments will depend on what the other customized
  macros do. For example if |opp=\foo| is used and the custom |\foo| inserts a
  |+| when taking the opposite of a negative number, then the other custom
  macros for arithmetic (and the |\foo| macro itself) must be able to handle
  arguments starting optionally with such a |+|.%
}
must be \emph{f}-expandable,
except for the computation of factorials, which only has to be
\emph{x}-expandable.%
%
\footnote{Prior to |1.4|, only \emph{x}-expandability was required. It is easy however to use an |\expanded| based wrapper to convert  \emph{x}-expandable macros into  \emph{f}-expandable ones.}
%

Macro |\bnumsetup| can be used multiple times in the same document, thus
allowing to switch math engines or to remap operators to some other arithmetic
macros of the same math engine. Its effect obeys the local scope.%
%
\footnote{The effect is global if under |\xintglobaldefstrue| setting.}

The hexadecimal input via the |"| prefix is converted internally to decimal
notation using |\xintHexToDec| from package \xintbinhexname, and customization
is possible via redefinition of |\bnumhextodec|, whose default is to be an
alias to |\xintHexToDec|.

The final conversion back to hexadecimal done by |\evaltohex| is handled by
|\bnumprintonetohex| which defaults to |\xintDecToHex|.

These two steps can thus be customized as will. But the loading of package
\xintbinhexname can not be canceled.

\section{\csh{bnumprintone}, \csh{bnumprintonetohex}, \csh{bnumprintonesep}}

The computed values are printed one by one, separated by a comma and a space
(this is customizable as |\bnumprintonesep|), and each value being handed over
to |\bnumprintone|. By default this does nothing else than producing its
argument as is, it can be redefined at will (perhaps using macros such as in
\autoref{sec:printing} to handle the case of very long numbers).

There is also |\bnumprintonetohex| which is used by |\evaltohex| (this is its
sole difference from |\bnumeval|).  Its default definition makes it an alias
to |\xintDecToHex| from package \xintbinhexname.

\section{Example of customization: let the syntax handle fractions!}

I will show how to transform |\bnumeval| into a calculator with fractions! We
will use the \xintfracname macros, but coerce them into always producing
fractions in lowest terms (except for powers). For optimization we use
the |[0]| post-fix which speeds-up the input parsing by the \xintfracname
macros. We remove it on output via a custom |\bnumprintone|.

Note that the |/| operator is associated to |divround| key but of course here
the used macro will simply do an exact division of fractions, not a rounded-to-%
an integer division. This is the whole point of using a macro of our own
choosing!
\begin{verbatim}
\usepackage{xintfrac}
\newcommand\myIrrAdd[2]{\xintIrr{\xintAdd{#1}{#2}}[0]}
\newcommand\myIrrSub[2]{\xintIrr{\xintSub{#1}{#2}}[0]}
\newcommand\myIrrMul[2]{\xintIrr{\xintMul{#1}{#2}}[0]}
\newcommand\myDiv[2]{\xintIrr{\xintDiv{#1}{#2}}[0]}
\newcommand\myDivFloor[2]{\xintDivFloor{#1}{#2}[0]}
\newcommand\myMod[2]{\xintIrr{\xintMod{#1}{#2}}[0]}
\newcommand\myPow[2]{\xintPow{#1}{#2}}% will have trailing [0]
\newcommand\myFac[1]{\xintFac{#1}}%     produces trailing [0]
\bnumsetup{add=\myIrrAdd, sub=\myIrrSub, mul=\myIrrMul,
           divround=\myDiv, div=\myDivFloor,
           mod=\myMod, pow=\myPow, fac=\myFac}%
% % if any operation happened, the result is already irreducible
% % (except power of non-irreducible) so this is overhead:
% \let\bnumprintone\xintIrr
% % but it is safe way to get rid of the trailing [0]
% % else we have to take care of case with no [0] because of no operations
% % well let's do it:
\makeatletter
\def\myPrintOne#1{\@myPrintOne#1[0]\relax}
\def\@myPrintOne#1[0]#2\relax{#1}
\let\bnumprintone\myPrintOne
\makeatother
\end{verbatim}
\begingroup
\newcommand\myIrrAdd[2]{\xintIrr{\xintAdd{#1}{#2}}[0]}
\newcommand\myIrrSub[2]{\xintIrr{\xintSub{#1}{#2}}[0]}
\newcommand\myIrrMul[2]{\xintIrr{\xintMul{#1}{#2}}[0]}
\newcommand\myDiv[2]{\xintIrr{\xintDiv{#1}{#2}}[0]}
\newcommand\myDivFloor[2]{\xintDivFloor{#1}{#2}[0]}
\newcommand\myMod[2]{\xintIrr{\xintMod{#1}{#2}}[0]}
\newcommand\myPow[2]{\xintPow{#1}{#2}}% will have trailing [0]
\newcommand\myFac[1]{\xintFac{#1}}%     produces trailing [0]
\bnumsetup{add=\myIrrAdd, sub=\myIrrSub, mul=\myIrrMul, divround=\myDiv,
  div=\myDivFloor, mod=\myMod, pow=\myPow, fac=\myFac}%
\makeatletter
\def\myPrintOne#1{\@myPrintOne#1[0]\relax}%
\def\@myPrintOne#1[0]#2\relax{#1}%
\let\bnumprintone\myPrintOne
\makeatother
\bneshoweval{1000000*(1/100+1/2^7-20/5^4)/(1/3-5/7+9/11)^2}
\bneshoweval{(1-1/2)(1-1/3)(1-1/4)(1-1/5)(1-1/6)(1-1/7)}
\bneshoweval{(1-1/3+1/9-1/27-1/81+1/243-1/729+1/2187)^5}
\bneshoweval{(1+1/10)^10 /: (1-1/10)^10}
\bneshoweval{2^-3^4}
\endgroup

This last example is computed differently than it would be with \xintexprname
|1.4f| because \bnumname |1.5| already applies right associativity to powers,
whereas \xintexprname |1.4f| still applies left associativity.

Note also that the above changes break |\evaltohex| whose output routine uses
by default |\xintDecToHex| which will choke on fractional input. However it is
not difficult to write a macro applying separately to numerator and
denominator.

Computations with fractions quickly give birth to big results, see
\autoref{sec:printing} on how to modify |\bnumprintone| to coerce \TeX\ into
wrapping numbers too long for the available width.

\section{Differences from \csh{numexpr}}
\label{sec:differences}

Apart from the extension to big integers (i.e. exceeding the \TeX{} limit at
\number"7FFFFFFF), and the added operators, there are a number of important
differences between |\bnumexpr| and |\numexpr|:

\begin{enumerate}
\item contrarily to |\numexpr|, the |\bnumexpr| parser stops only after having
  found (and swallowed) a mandatory ending |\relax| token (it can arise from expansion),

\item in particular note that spaces between digits do not stop |\bnumexpr|, in contrast
  with |\numexpr|:

  |\the\numexpr 3 5+79\relax| expands (in one step) to \expandafter|\the\numexpr 3 5+79\relax|

  |\thebnumexpr 3 5+79\relax| expands (in two steps) to
  \expandafter\expandafter\expandafter|\thebnumexpr 3 5+79\relax|

\item with |\edef\myVar{\thebnumexpr1+2\relax}|, the computation is of course
  done at time of the |\edef|. But one is also allowed to do
  |\edef\myVar{\bnumexpr1+2\relax}| which prepares |\myVar| as a macro which
  can be inserted in other \bnumname expressions and behave there as a
  self-contained pre-computed unit triggering tacit multiplication, or be
  typeset directly if inserted in the typesetting stream.%
%
\footnote{Prior to |1.4|, one would have had to use |\bnethe\myVar| for
  typesetting, or |\bnumeval{\myVar}|.}
%
  There is no analog with |\numexpr| as |\edef\myVar{\numexpr1+2\relax}| does
  not pre-compute anything and furthermore |\the\numexpr2\myVar\relax| in
  typesetting flow then triggers the |You can't use `\numexpr' in horizontal
  mode| error.

\item expressions may be comma separated. On input, spaces are ignored,
  and on output the values are comma separated with a space after
  each comma,

\item |\bnumexpr -(1+1)\relax| is legal contrarily to |\numexpr -(1+1)\relax|
  which raises an error,

\item |\numexpr 2\cnta\relax| is illegal (with |\cnta| a |\count|-variable.)
  But |\bnumexpr 2\cnta\relax| is perfectly legal and will do the tacit
  multiplication,

\item more generally, tacit multiplication applies in front of parenthesized
  sub-expressions, or sub |\bnumexpr...\relax| (or |\numexpr...\relax|), or
  also after parentheses in front of numbers,

\item the underscore |_| is accepted within the digits composing a number and
  is silently ignored by |\bnumexpr|.

\end{enumerate}

As hinted above |\bnumexpr...\relax| differs from |\thebnumexpr...\relax| as
the latter expands to explicit digit tokens, but the former expands to a
private self-contained format which can serve as sub-unit in other
expressions, or be used inside |\edef|.  Since |1.4| the former idiom can also
be inserted directly inside the typesetting stream, or be written out to an
external file where it will expand to some control sequences, braces, and
character tokens, all with their standard catcodes. 

One can use |\numexpr...\relax| as a sub-unit in |\bnumexpr...\relax| but the
reverse does not apply: it would either cause an error or an
anticipated end to the |\numexpr| which will think having hit a |\relax|.

An important thing to keep in mind is that if one has a calculation
whose result is a small integer, acceptable by \TeX{} in |\ifnum| or
count assignments, this integer produced by |\thebnumexpr| is not
self-delimiting, contrarily to a |\numexpr...\relax| construct: the
situation is exactly as with a |\the\numexpr...\relax|, thus one may
need to terminate the number to avoid premature expansion of following
tokens; for example with the |\space| control sequence. When using |\bnumeval{...}|
syntax as in
\begin{verbatim}
\ifnum\bnumeval{...}
...
\fi
\end{verbatim}
the end of line will (under the normal \LaTeX\ configuration) insert a
terminating space token. Again, here |\bnumeval{...}| must produce an integer
acceptable to \TeX, i.e. at most {\number"7FFFFFFF} in absolute value.

\section{Printing big numbers}
\label{sec:printing}
\LaTeX{} will not split long numbers at the end of lines. I personally
often use helper macros (not in the package) of the following type:

\begin{verbatim}
\def\allowsplits #1{\ifx #1\relax \else #1\hskip 0pt plus 1pt\relax
                    \expandafter\allowsplits\fi}%
\def\printnumber #1{\expandafter\allowsplits \romannumeral-`0#1\relax }%
% \printnumber thus first ``fully'' expands its argument.
\end{verbatim}

\def\allowsplits #1{\ifx #1\relax \else #1\hskip 0pt plus 1pt\relax
                    \expandafter\allowsplits\fi}%
\def\printnumber #1{\expandafter\allowsplits \romannumeral-`0#1\relax }%

\noindent|\thebnumexpr 1000!\relax=|\phantom{0}{\color{digitscolor}\printnumber{\thebnumexpr 1000!\relax}}\par

\section{Expression syntax and its customizability}

The implemented syntax is the expected one with infix operators and
parentheses, the recognized operators being |+|, |-|, |*|, |/| (rounded
division), |^| (power), |**| (power), |//| (by default floored division), |/:|
(the associated modulo) and |!| (factorial). One can input hexadecimal numbers
as in \TeX\ syntax for number assignments, i.e. using a |"| prefix and only
uppercase letters |ABCDEF|.

Different computations may be separated by commas. The whole expression is
handled token by token, any component (digit, operator, parentheses... even
the ending |\relax|) may arise on the spot from macro expansions.  The
underscore |_| can be used to separate digits in long numbers, for readability
of the input.

The precedence rules are as expected and detailed in the next section.
Operators on the same level of precedence (like |*|, |/|, |//|, |/:|) behave
in a left-associative way, and these examples behave as e.g.\@ with Python
analogous operators:\newline \bneshoweval{100//3*4, 100*4//3, 100/:3*4,
  100*4/:3, 100//3/:5}

At |1.5| a change was made to the power operators which became
right-associative. Again, this matches the behaviour e.g.\@ of Python:%
%
\footnote{It had been announced at \xintexprname \texttt{1.4} that probably in
  future power operator would become right-associative, so we experiment it
  here in \bnumname in advance.}%
%
\newline
\bneshoweval{2^3^4, 2^(3^4)}

It is possible to customize completely the behaviour of the parser, in two ways:
\begin{itemize}
\item via |\bnumsetup| which has a simple interface to replace the macros
  associated to the operators |+|, |-|, |*|, |/|, |//|, |/:|, |**| and |^| by
  custom macros,
\item or even more completely via |\bnumdefinfix| and |\bnumdefpostfix| which
  allow to add new operators to the syntax! (or overwrite existing ones...)
\end{itemize}

\section{Precedences}

The parser implements precedence rules based on concepts which are summarized
below. I am providing them for users who will use the customizing macros.
\begin{itemize}
\item an infix operator has two associated precedence levels, say |L| and
  |R|,
\item the parser proceeds from left to right, pausing each time it has found a
  new number and an operator following it,
\item the parser compares the left-precedence |L| of the new found operator to
  the right-precedence |R_last| of the last delayed operation (which already
  has one argument and would like to know if it can use the new found one): if
  |L| is at most equal to it, the delayed operation is now executed, else the
  new-found operation is kept around to be executed first, once it will have
  gathered its arguments, of which only one is known at this stage.
\end{itemize}

Although there is thus internally all the needed room for sophistication, the
implemented table of precedences simply puts all of multiplication and
division related operations at the same level, which means that left
associativity will apply with these operators. I could see that Python behaves
the same way for its analogous operators.

Here is the default table of precedences as implemented by the package:
\begin{center}
\DeleteShortVerb{\|}%
  \begin{tabular}{|c|c|c|}
    \multicolumn3{c}{Table of precedences}\\\hline
\MakeShortVerb{\|}%
    operator& left &right\\\hline
    |+|,|-|& 12&12\\
    |*|,|/|,|//|,|/:|& 14&14\\
    tacit |*|&16&14\\
    |**|, |^|& 18&17\\
    |!| &20&n/a\\\hline
  \end{tabular}
\end{center}
Tacit multiplication applies in front of parentheses, and after them, also in
front of count variables or registers. As shown in the table it has an
elevated precedence compared to multiplication explicitly induced by |*|, so
|100/4(9)| is computed as |100/36| and not as |25*9|:\newline
\bneshoweval{100/4(9), (100/4)9, 1000 // (100/4) 9 (1+1) * 13}
More generally |A/B(C)(D)(E)*F| will compute |(A/(B*C*D*E))*F|.%
%
\footnote{The |B(C)(D)(E)| product will be computed as |B*(C*(D*E))| because
  the right-precedence of tacit multiplication is |14| but its left-precedence
  is |16|, creating right associativity.  As the underlying mathematical
  operation is associative this is irrelevant to final result.}

The unary |-|, as prefix, has a special behaviour: after an infix operator it
will acquire a right-precedence which is the minimum of |12| (i.e.\@ the
precedence of addition and subtraction) and of the right-precedence of the
infix operator.  For example |2^-3^4| will be parsed as |2^(-(3^4))|, raising
an error because the parser is by default integer only, but see the section
about |\bnumsetup| which explains how to let |\bnumeval| computes fractions!


\section{\csh{bnumdefinfix}}

It is possible to define infix binary operators of one's own choosing.%
%
\footnote{The effect of |\bnumdefinfix| is global if under |\xintglobaldefstrue| setting.}
%
The syntax is 
\centeredline{|\bnumdefinfix|\marg{operator}\marg{\string\macro}\marg{L-prec}\marg{R-prec}}
\begin{description}
\item[\marg{operator}] The characters for the operator, they may be letters or
  non-letters, and must not be active or among the special characters |\|,
  \textcolor{smallverbcolor}{\{}, \textcolor{smallverbcolor}{\}}, |#| and
  {\catcode`\% 12 |%|}.  Also,
    spaces will be removed.%
%
\footnote{The |_| can be used, but not as first character of the operator, as
  it would be mis-construed on usage as part of the previous number, and
  ignored as such.}%
%
$^{,}$%
%
{\catcode`# 12 \footnote{It is actually possible to use |#| as an operator
    name or a character in such a name but the definition with |\bnumdefinfix|
    must then be done either with \expandafter|\string\string#| or
    |####|...}}%
$^{,}$%
%
\footnote{Active characters must be prefixed by |\string| both at the time of
  the definition and at the time of use. It is probably better to use a toggle which
  will turn off the activity, both at time of definition and at time of use.}%
\item[\marg{\string\macro}] The expandable macro (expecting two mandatory
  arguments) which is to assign to the infix operator. This macro must be
  \emph{f}-expandable. Also it must (if the default package configuration is
  not modified for the core operators) produce integers in the ``strict''
  format which is expected by the \xintcorename macros for arithmetic: no
  leading zeros, at most one minus sign, no plus sign, no spaces.
\item[\marg{L-prec}] An integer, minimal |4|, maximal |22|, which governs
  the left-\hskip0pt precedence of the infix operator.
\item[\marg{R-prec}] An integer, minimal |4|, maximal |22|, which governs the
  right-\hskip0pt precedence of the infix operator.
\end{description}
Generally, the two precedences are set to the same value.

Once a multi-character operator is defined, the first characters of its name
can be used if no ambiguity.  In case of ambiguity, it is the earliest defined
shortcut which prevails, except for the full name. So for example if |$abc|
operator is defined, and |$ab| is defined next, then |$| and |$a| will still
serve as shortcuts to the original |$abc|, but |$ab| will refer to the newly
defined operator.

Fully qualified names are never ambiguous, and a shortcut once defined will
change meaning under only these two possibilities:
\begin{itemize}
\item it is re-defined as the full name of a new operator,
\item the original operator to which the shortcut refers is defined again; then 
  the shortcut is automatically updated to point to the new meaning.
\end{itemize}
\begin{verbatim}
\def\equals#1#2{\ifnum\pdfstrcmp{#1}{#2}=0 \expandafter1\else
                                             \expandafter0\fi}
% or:
\def\equals#1#2{\expanded{\ifnum\pdfstrcmp{#1}{#2}=0 1\else0\fi}}
\def\differ#1#2{\expanded{\ifnum\pdfstrcmp{#1}{#2}=0 0\else1\fi}}
\bnumdefinfix{==}{\equals}{10}{10}
\bnumdefinfix{!=}{\differ}{10}{10}
\bnumdefinfix{times}{\xintiiMul}{14}{14}
\bnumdefinfix{++}{\xintiiAdd}{19}{19}
\end{verbatim}
%\def\equals#1#2{\ifnum\pdfstrcmp{#1}{#2}=0 \expandafter1\else\expandafter0\fi}
\def\equals#1#2{\expanded{\ifnum\pdfstrcmp{#1}{#2}=0 1\else0\fi}}
\def\differ#1#2{\expanded{\ifnum\pdfstrcmp{#1}{#2}=0 0\else1\fi}}
\bnumdefinfix{==}{\equals}{10}{10}
\bnumdefinfix{!=}{\differ}{10}{10}
\bnumdefinfix{times}{\xintiiMul}{14}{14}
\bnumdefinfix{++}{\xintiiAdd}{19}{19}
%
\bneshoweval{2 + 3! = 5, 2 + (3!) == 8} Notice in the |2+3! = 5| example that
the existence of |!=| prevails on applying the factorial, so this is test
whether |2+3| and |5| differ; it is not a matter of precedence here, but of
input parsing ignoring spaces. And |2+3! == 8| would create an error as after
having found the |!=| operator and now expecting a digit (as there is no |!==|
operator) the parser would find an unexpected |=| and report an error. Hence
the usage of parentheses in the input.%
%
\footnote{As |!=| is indeed defined out-of-the-box in the \xintexprname
  syntax, the |3! == 8| issue applies with |\xinteval| and perhaps I should
  add it to the user documentation, as warning.}\newline
\bneshoweval{2^5 == 4 times 8, 11 t 14}
\bneshoweval{100 ++ -10 ^ 3, (100 - 10)^3, 2 ** 5 ++ 3, 2^(5+3)}


\section{\csh{bnumdefpostfix}}

It is possible to define postfix unary operators of one's own choosing.%
%
\footnote{The effect of |\bnumdefpostfix| is global if under |\xintglobaldefstrue| setting.}
%
The syntax is 
\centeredline{|\bnumdefpostfix|\marg{operator}\marg{\string\macro}\marg{L-prec}}
\begin{description}
\item[\marg{operator}] The characters for the operator name: same
  conditions as for |\bnumdefinfix|. Postfix and infix operators share the
  same name-space, regarding abbreviated names.
\item[\marg{\string\macro}] The one argument expandable macro to assign to
  the postfix operator. This macro only needs to be \emph{x}-expandable.
\item[\marg{L-prec}] An integer, minimal |4|, maximal |22|, which governs the
  left-\hskip0pt precedence of the infix operator.
\end{description}
Examples below which use the maximal precedence are typical of what is
expected of a ``function'' (and I even used |.len()| notation with parentheses
in one example, the parentheses are part of the postfix operator name).  And
indeed such postfix operators are thus a way to implement functions in
disguise, circumventing the fact that the \bnumname parser will never be
extended to work with functional syntax (for this, see \xintexprname).  With
the convention (followed in some examples) that such postfix operators start
with a full stop, but never contain another one, we can chain simply by using
concatenation (no need for parentheses), as there will be no
ambiguity.
\begin{verbatim}
\usepackage{xint}% for \xintiiSum, \xintiiSqrt
\def\myRev#1{\xintNum{\xintReverseOrder{#1}}}% reverse and trim leading zeros
\bnumdefpostfix{$}{\myRev}{22}%     the $ will have top precedence
\bnumdefpostfix{:}{\myRev}{4}%      the : will have lowest precedence
\bnumdefpostfix{::}{\xintiiSqr}{4}% the :: is a completely different operator
\bnumdefpostfix{.len()}{\xintLength}{22}% () for fun but a single . will be enough!
\bnumdefpostfix{.sumdigits}{\xintiiSum}{22}% .s will abbreviate
\bnumdefpostfix{.sqrt}{\xintiiSqrt}{22}%     .sq will be unambiguous (but confusing)
\bnumdefpostfix{.rep}{\xintReplicate3}{22}%  .r will be unambiguous
\end{verbatim}
\def\myRev#1{\xintNum{\xintReverseOrder{#1}}}% reverse digits and trim leading zeros
\bnumdefpostfix{$}{\myRev}{22}% $
\bnumdefpostfix{:}{\myRev}{4}%
\bnumdefpostfix{::}{\xintiiSqr}{4}%
\bnumdefpostfix{.len()}{\xintLength}{22}%
\bnumdefpostfix{.sumdigits}{\xintiiSum}{22}%
\bnumdefpostfix{.sqrt}{\xintiiSqrt}{22}%
\bnumdefpostfix{.rep}{\xintReplicate3}{22}%
\bneshoweval{(2^31).len(), (2^31)., 2^31$, 2^31:, (2^31)$}
\bneshoweval{(2^31).sqrt, 100000000.sq.sq}
\bneshoweval{(2^31).sumdigits, 123456789.s, 123456789.s.s, 123456789.s.s.s}
\bneshoweval{10^10+10000+2000+300+40+5:}
\bneshoweval{1+2+3+4+5+6+7+8+9+10 :: +1 :: *2 :: :: :}
\bneshoweval{123456789.r}\par\vspace*{-\baselineskip}% Toujours LaTeX avec ses
                                % espaces verticaux TROP GRANDS DEPUIS 30 ANS
\begin{verbatim}
\bnumdefpostfix{.rep}{\xintReplicate5}{22}% .rep modified --> .r too
\end{verbatim}
\bnumdefpostfix{.rep}{\xintReplicate5}{22}%
\bneshoweval{123456789.r}

\clearpage

\section{Readme}
\begingroup
\makeatletter\def\x{\baselineskip10pt
                    \ttfamily\settowidth\dimen@{X}%
                    %\parindent \dimexpr.5\linewidth-34\dimen@\relax
                    \parindent\z@
                    \let\do\do@noligs\verbatim@nolig@list
                    \let\do\@makeother\dospecials
                    \makestarlowast
                    %\def\par{\leavevmode\null\@@par\penalty\interlinepenalty}%
                    \def\par{\leavevmode\null\@@par\pagebreak[1]}%
                    \@vobeyspaces\obeylines
                    \noindent\kern\parindent\input README.md
\endgroup }\x

\clearpage

\section{Changes}
%\small
\begin{description}
\input bnumexprchanges.tex
\end{description}

\StopEventually{\end{document}\endinput}

\newgeometry{hscale=0.75,vscale=0.75}% ATTENTION \newgeometry fait
                                % un reset de vscale si on ne le
                                % précise pas ici !!!

\MakePercentIgnore
%
% \catcode`\<=0 \catcode`\>=11 \catcode`\*=11 \catcode`\/=11
% \let</dtx>\relax
% \def<*package>{\catcode`\<=12 \catcode`\>=12 \catcode`\*=12 \catcode`\/=12 }
%</dtx>
%<*package>
%
% \section{\bnumname implementation}
% \label{sec:bnumexprcode}
% \etocdefaultlines
% \localtableofcontents
% \etocmarkbothnouc{\bnumnameuserdoc \hyperref[sec:bnumexprcode]{implementation}}
%
% I transferred mid-May 2021 from \xintexprname its |\expanded| based
% infra-structure from its own |1.4| release of January 2020 and bumped
% version to |1.4|. Also I added support for hexadecimal input and output, via
% unconditional loading of \xintbinhexname.
%
% A few comments added here at |1.4a|:
% \begin{itemize}
% \item It looked a bit costly and probably would have been mostly useless to
%   end users to integrate in \bnumname support for nested structures via
%   square brackets [, ], which is in \xintexprname since its January 2020
%   |1.4| release.  But some of the related architecture remains here; we
%   could make some gains probably but diverging from upstream code
%   would make maintenance a nightmare.
% \item Formerly, the |\csname...\endcsname| encapsulation technique had the
%   after-effect to allow the macros supporting the infix operators to be only
%   \emph{x}-expandable.  At |1.4|, I could have still allowed support macros
%   being only \emph{x}-expandable, but, keeping in sync with upstream, I have
%   used only a |\romannumeral| trigger and did not insert an |\expanded|, so
%   now the support macros must be \emph{f}-expandable.  The |1.4a| release
%   fixes the related user documentation of |\bnumsetup| which was not
%   updated at |1.4|.  The support macro for the factorial however needs only
%   be \emph{x}-expandable.
% \item Also, I simply do not understand why the legacy (|1.2e|) user documentation
%   said that the support macros were supposed to \emph{f}-expand their arguments, as
%   they are used only with arguments being explicit digit tokens (and
%   optional minus sign).
% \item The |\bnumexpr\relax| syntax creating an empty ople is by itself now
%   legal, and can be injected (comma separated) in an expression, keeping it
%   invariant, however |\bnumeval{}| ends in a |File ended while scanning use
%   of \BNE_print_c| error because |\BNEprint| makes the tacit requirement
%   that the 1D ople to output has at least one item.
% \end{itemize}
%
% At |1.5|, right-associativity is implemented for powers in anticipation of
% upstream, and the customizability and extendibility of the package is made
% total via added |\bnumdefinfix| and |\bnumdefpostfix|.
%
% \subsection{Package identification}\leavevmode
%    \begin{macrocode}
\NeedsTeXFormat{LaTeX2e}%
\ProvidesPackage{bnumexpr}[2021/05/17 v1.5 Expressions with big integers (JFB)]%
%    \end{macrocode}
% \subsection{Load unconditionally xintbinhex}
% Newly done at |1.4|. Formerly, \bnumname had no dependency if loaded
% with option |custom|. But for |1.4| release I have decided to add
% unconditional support for hexadecimal notation.
%
% Let's require the most recent \xintname date at time of
% writing. We should check for availability of |\expanded| but well.
%
%    \begin{macrocode}
\RequirePackage{xintbinhex}[2021/05/10]%
%    \end{macrocode}
% \subsection{Package options}\leavevmode
%    \begin{macrocode}
\def\BNEtmpa {0}%
\DeclareOption {custom}{\def\BNEtmpa {1}}%
\ProcessOptions\relax
%    \end{macrocode}
% \subsection{\csh{bnumsetup} and conditional loading of \xintcorename}
% The keys should have been |Add|, |Sub|, \dots, not |add|, |sub|, \dots, so
% internally macros |\BNE_Op_Add| etc\dots\ macro names would be used, but well,
% let's simply live with this.
%
% |\bnumsetup| replaces at |1.5| deprecated |\bnumexprsetup| which is kept as an alias.
%    \begin{macrocode}
{\catcode`! 3 \catcode`_ 11 %
  \gdef\bnumsetup #1{\BNE_parsekeys #1,=!,}%
  \gdef\BNE_parsekeys #1=#2#3,%
  {%
    \ifx!#2\expandafter\BNE_parsedone\fi
  \XINT_global
    \expandafter
    \let\csname BNE_Op_\xint_zapspaces #1 \xint_gobble_i\endcsname%
    =#2%
  \ifxintverbose
    \PackageInfo{bnumexpr}{assigned
    \ifxintglobaldefs globally \fi
     \string#2 to \xint_zapspaces #1 \xint_gobble_i\MessageBreak
%    \end{macrocode}
% Workaround the space inserted by |\on@line|.
%    \begin{macrocode}
     \expandafter\xint_firstofone}%
  \fi
  \BNE_parsekeys
  }%
  \gdef\BNE_parsedone #1\BNE_parsekeys {}%
}%
\let\bnumexprsetup\bnumsetup
\if0\BNEtmpa\expandafter\@secondoftwo\fi
\@gobble{%
    \RequirePackage{xintcore}[2021/05/10]%
    \bnumsetup{add=\xintiiAdd, sub=\xintiiSub, mul=\xintiiMul,
                   divround=\xintiiDivRound, div=\xintiiDivFloor,
                   mod=\xintiiMod, pow=\xintiiPow, fac=\xintiiFac,
                   opp=\xintiiOpp}%
}%
%    \end{macrocode}
% \subsection{Activate usual \xintname catcodes for code source}\leavevmode
%    \begin{macrocode}
\edef\BNErestorecatcodes{\XINTrestorecatcodes}%
\XINTsetcatcodes%
%    \end{macrocode}
% Strangely those three are not defined in xintkernel.sty, but only in xint.sty
%    \begin{macrocode}
\long\def\xint_firstofthree  #1#2#3{#1}%
\long\def\xint_secondofthree #1#2#3{#2}%
\long\def\xint_thirdofthree  #1#2#3{#3}%
%    \end{macrocode}
% For the mechanism of |\bnumdefinfix| we need [1] precedence levels to be
% available as |\chardef|'s. \xintkernelname already provides 0-10, 12, 14,
% 16, 18, 20, 22.  Admittedly they could be created only dynamically, and then
% I would not have to set a cap at 22, but well, that's already a large
% supported range for a functionality nobody will use, as nobody probably uses
% the package to start with.
%
% .. [1] left levels need to be represented by one token; right levels are
% hard-coded into |checkp_<op>| macros and could have been there explicit
% digit tokens but we will use the |\xint_c_...| |\char|-tokens.
%    \begin{macrocode}
\chardef\xint_c_xi   11
\chardef\xint_c_xiii 13
\chardef\xint_c_xv   15
\chardef\xint_c_xvii 17
\chardef\xint_c_xix  19
\chardef\xint_c_xxi  21
%    \end{macrocode}
% \subsection{\csh{bnumexpr}, \csh{thebnumexpr}, \csh{bnethe}, \csh{bnumeval}}\leavevmode
%    \begin{macrocode}
\def\XINTfstop {\noexpand\XINTfstop}%
\def\bnumexpr  {\romannumeral0\bnumexpro}%
\def\bnumexpro {\expandafter\BNE_wrap\romannumeral0\bnebareeval }%
\def\BNE_wrap  {\XINTfstop\BNEprint.}%
\def\bnumeval #1%
   {\expanded\expandafter\BNEprint\expandafter.\romannumeral0\bnebareeval#1\relax}%
\def\evaltohex #1%
   {\expanded\expandafter\BNEprinthex\expandafter.\romannumeral0\bnebareeval#1\relax}%
\def\thebnumexpr
   {\expanded\expandafter\BNEprint\expandafter.\romannumeral0\bnebareeval}%
\def\bnebareeval{\BNE_start}%
\def\bnethe#1{\expanded\expandafter\xint_gobble_i\romannumeral`&&@#1}%
\protected\def\BNEprint.#1{{\BNE_print#1.}}%
\def\BNE_print#1{\bnumprintone{#1}\expandafter\BNE_print_a\string}%
\def\BNE_print_a#1{\unless\if#1.\expandafter\BNE_print_b\fi}%
\def\BNE_print_b
   {\expandafter\BNE_print_c\expandafter{\expandafter\xint_gobble_i\string}}%
\def\BNE_print_c#1{\bnumprintonesep\bnumprintone{#1}\expandafter\BNE_print_a\string}%
\protected\def\BNEprinthex.#1{{\BNE_printhex#1.}}%
\def\BNE_printhex#1{\bnumprintonetohex{#1}\expandafter\BNE_printhex_a\string}%
\def\BNE_printhex_a#1{\unless\if#1.\expandafter\BNE_printhex_b\fi}%
\def\BNE_printhex_b
   {\expandafter\BNE_printhex_c\expandafter{\expandafter\xint_gobble_i\string}}%
\def\BNE_printhex_c#1{\bnumprintonesep\bnumprintonetohex{#1}\expandafter\BNE_printhex_a\string}%
\let\bnumprintone\xint_firstofone
\let\bnumprintonetohex\xintDecToHex
\def\bnumprintonesep{, }%
%    \end{macrocode}
% \subsection{\csh{BNE_getnext}}
% The upstream |\BNE_put_op_first| has a string of included |\expandafter|,
% which was imported here at |1.4| and |1.4a| but they serve nothing in our
% context. Removed this useless overhead at |1.5|.
%    \begin{macrocode}
\def\BNE_getnext #1%
{%
    \expandafter\BNE_put_op_first\romannumeral`&&@%
    \expandafter\BNE_getnext_a\romannumeral`&&@#1%
}%
\def\BNE_put_op_first #1#2#3{#2#3{#1}}%
\def\BNE_getnext_a #1%
{%
    \ifx\relax #1\xint_dothis\BNE_foundprematureend\fi
    \ifx\XINTfstop#1\xint_dothis\BNE_subexpr\fi
    \ifcat\relax#1\xint_dothis\BNE_countetc\fi
    \xint_orthat{}\BNE_getnextfork #1%
}%
\def\BNE_foundprematureend\BNE_getnextfork #1{{}\xint_c_\relax}%
\def\BNE_subexpr #1.#2%
{%
    \expanded{\unexpanded{{#2}}\expandafter}\romannumeral`&&@\BNE_getop
}%
\def\BNE_countetc\BNE_getnextfork#1%
{%
    \if0\ifx\count#11\fi
        \ifx\dimen#11\fi
        \ifx\numexpr#11\fi
        \ifx\dimexpr#11\fi
        \ifx\skip#11\fi
        \ifx\glueexpr#11\fi
        \ifx\fontdimen#11\fi
        \ifx\ht#11\fi
        \ifx\dp#11\fi
        \ifx\wd#11\fi
        \ifx\fontcharht#11\fi
        \ifx\fontcharwd#11\fi
        \ifx\fontchardp#11\fi
        \ifx\fontcharic#11\fi 0\expandafter\BNE_fetch_as_number\fi
   \expandafter\BNE_getnext_a\number #1%
}%
\def\BNE_fetch_as_number
    \expandafter\BNE_getnext_a\number #1%
{%
    \expanded{{{\number#1}}\expandafter}\romannumeral`&&@\BNE_getop
}%
%    \end{macrocode}
% In the case of hitting a |(|, previous release inserted directly a
% |\BNE_oparen|.  But the expansion architecture imported from upstream
% |\xintiiexpr| has been refactored, and the |..._oparen| meaning and
% usage evolved.  We stick with |{}\xint_c_ii^v (| from upstream.
%    \begin{macrocode}
\def\BNE_getnextfork #1{%
    \if#1+\xint_dothis \BNE_getnext_a \fi
    \if#1-\xint_dothis {{}{}-}\fi
    \if#1(\xint_dothis {{}\xint_c_ii^v (}\fi
    \xint_orthat {\BNE_scan_number #1}%
}%
%    \end{macrocode}
% \subsection{Parsing an integer in decimal or hexadecimal notation}\leavevmode
%    \begin{macrocode}
\def\BNE_scan_number #1%
{%
    \if "#1\xint_dothis \BNE_scanhex\fi
    \ifnum \xint_c_ix<1\string#1 \xint_dothis \BNE_startint\fi
    \xint_orthat \BNE_notadigit #1%
}%
%    \end{macrocode}
% If user employs |\bnumdefinfix| with |\string#|, and then tries |100##3|,
% the first |#| will be interpreted as operator (assuming no
% operator starting with |##| has actually been defined) and the error
% "message" (which is not using |\message| or a |\write|)
% will then be
% \centeredline{|Digit? (got `##')|}
% because the parser is actually looking for a digit but finds the second |#|,
% and TeX displays it doubled. This is doubly
% confusing, but well, let's not dwell on that.
%    \begin{macrocode}
\def\BNE_notadigit#1%
{%
    \expandafter\BNE_scan_number
    \romannumeral`&&@\XINT_expandableerror{Digit? (got `#1'). Hit I<RET><digit>}%
}%
\def\BNE_startint #1%
{%
    \if #10\expandafter\BNE_gobz_a\else\expandafter\BNE_scanint_a\fi #1%
}%
\def\BNE_scanint_a #1#2%
    {\expanded\bgroup{{\iffalse}}\fi #1%
     \expandafter\BNE_scanint_main\romannumeral`&&@#2}%
\def\BNE_gobz_a #1#2%
    {\expanded\bgroup{{\iffalse}}\fi
     \expandafter\BNE_gobz_scanint_main\romannumeral`&&@#2}%
\def\BNE_scanint_main #1%
{%
    \ifcat \relax #1\expandafter\BNE_scanint_hit_cs \fi
    \ifnum\xint_c_ix<1\string#1 \else\expandafter\BNE_scanint_next\fi
    #1\BNE_scanint_again
}%
\def\BNE_scanint_again #1%
{%
    \expandafter\BNE_scanint_main\romannumeral`&&@#1%
}%
%    \end{macrocode}
% Upstream (at |1.4f|) has |_getop| here, but let's jump directly to |BNE_getop_a|.
%    \begin{macrocode}
\def\BNE_scanint_hit_cs \ifnum#1\fi#2\BNE_scanint_again
{%
    \iffalse{{{\fi}}\expandafter}\romannumeral`&&@\BNE_getop_a#2%
}%
\def\BNE_scanint_next #1\BNE_scanint_again
{%
    \if    _#1\xint_dothis\BNE_scanint_again\fi
    \xint_orthat
    {\iffalse{{{\fi}}\expandafter}\romannumeral`&&@\BNE_getop_a#1}%
}%
\def\BNE_gobz_scanint_main #1%
{%
    \ifcat \relax #1\expandafter\BNE_gobz_scanint_hit_cs\fi
    \ifnum\xint_c_x<1\string#1 \else\expandafter\BNE_gobz_scanint_next\fi
    #1\BNE_scanint_again
}%
\def\BNE_gobz_scanint_again #1%
{%
    \expandafter\BNE_gobz_scanint_main\romannumeral`&&@#1%
}%
%    \end{macrocode}
% Upstream (at |1.4f|) has |_getop| here, but let's jump directly to |BNE_getop_a|.
%    \begin{macrocode}
\def\BNE_gobz_scanint_hit_cs\ifnum#1\fi#2\BNE_scanint_again
{%
    0\iffalse{{{\fi}}\expandafter}\romannumeral`&&@\BNE_getop_a#2%
}%
\def\BNE_gobz_scanint_next #1\BNE_scanint_again
{%
    \if    _#1\xint_dothis\BNE_gobz_scanint_again\fi
    \if    0#1\xint_dothis\BNE_gobz_scanint_again\fi
    \xint_orthat
    {0\iffalse{{{\fi}}\expandafter}\romannumeral`&&@\BNE_getop_a#1}%
}%
\def\BNE_hex_in #1.%
{%
    \expanded{{{\bnumhextodec{#1}}}\expandafter}\romannumeral`&&@\BNE_getop
}%
\let\bnumhextodec\xintHexToDec
%    \end{macrocode}
% Upstream (until |1.4f|) had a long-standing bug in its hexadecimal input,
% which was inherited here at |1.4|: the |\BNE_scanhex| triggered
% |\BNE_scanhex_a| which then grabbed an unexpanded token and used it as is
% in an |\ifcat|... this made syntax such as |"\foo| broken.  Fixed here at
% |1.5|.
%
% And there was a further long-standing bug in upstream (from |1.2m| to
% |1.4f|) about leading hexadecimal zeros not being trimmed. This was
% inherited here at |1.4|. Fixed also at |1.5|.
%    \begin{macrocode}
\def\BNE_scanhex #1#2% #1="
{%
    \expandafter\BNE_hex_in\expanded\bgroup
    \expandafter\BNE_scanhexgobz_a\romannumeral`&&@#2%
}%
\def\BNE_scanhexgobz_a #1%
{%
    \ifcat #1\relax0.\iffalse{\fi\expandafter}\expandafter\xint_gobble_i\fi
    \BNE_scanhexgobz_aa #1%
}%
\def\BNE_scanhexgobz_aa #1%
{%
    \if\ifnum`#1>`0
       \ifnum`#1>`9
       \ifnum`#1>`@
       \ifnum`#1>`F
       0\else1\fi\else0\fi\else1\fi\else0\fi 1%
       \xint_dothis\BNE_scanhex_b
    \fi
    \if 0#1\xint_dothis\BNE_scanhexgobz_bgob\fi
    \if _#1\xint_dothis\BNE_scanhexgobz_bgob\fi
    \if .#1\xint_dothis\BNE_scanhexgobz_toII\fi
    \xint_orthat
     {\XINT_expandableerror
        {HexDigit was expected but saw `#1'. Using 0, hit <RET>}%
      0.>;\iffalse{\fi}}%
    #1%
}%
\def\BNE_scanhexgobz_bgob #1#2%
{%
    \expandafter\BNE_scanhexgobz_a\romannumeral`&&@#2%
}%
\def\BNE_scanhex_a #1%
{%
    \ifcat #1\relax.\iffalse{\fi\expandafter}\expandafter\xint_gobble_i\fi
    \BNE_scanhex_aa #1%
}%
\def\BNE_scanhex_aa #1%
{%
    \if\ifnum`#1>`/
       \ifnum`#1>`9
       \ifnum`#1>`@
       \ifnum`#1>`F
       0\else1\fi\else0\fi\else1\fi\else0\fi 1%
       \expandafter\BNE_scanhex_b
    \else
       \if _#1\xint_dothis{\expandafter\BNE_scanhex_bgob}\fi
       \xint_orthat {.\iffalse{\fi\expandafter}}%
    \fi
    #1%
}%
\def\BNE_scanhex_b #1#2%
{%
    #1\expandafter\BNE_scanhex_a\romannumeral`&&@#2%
}%
\def\BNE_scanhex_bgob #1#2%
{%
    \expandafter\BNE_scanhex_a\romannumeral`&&@#2%
}%
%    \end{macrocode}
% \subsection{\csh{BNE_getop}}
% The upstream analog to |\BNE_getop_a| applies |\string| to |#1| in its
% thirdofthree branch before handing over to analog of |\BNE_scanop_a|, but I
% see no reason for doing it here (and I do have to check if upstream has
% any valid reason to do it). Removed. First branch was a |\BNE_foundend|,
% used only here, and expanding to |\xint_c_\relax|, let's move the
% |#1| (which will be |\relax|) last and simply insert |\xint_c_|.
%
% The |_scanop| macros have been refactored at upstream and here |1.5|.
%    \begin{macrocode}
\def\BNE_getop #1%
{%
    \expandafter\BNE_getop_a\romannumeral`&&@#1%
}%
\catcode`* 11
\def\BNE_getop_a #1%
{%
    \ifx   \relax #1\xint_dothis\xint_firstofthree\fi
    \ifcat \relax #1\xint_dothis\xint_secondofthree\fi
    \ifnum\xint_c_ix<1\string#1 \xint_dothis\xint_secondofthree\fi
    \if    (#1\xint_dothis      \xint_secondofthree\fi %)
    \xint_orthat \xint_thirdofthree
    \xint_c_
    {\BNE_prec_tacit *}%
    \BNE_scanop_a
    #1%
}%
\catcode`* 12
\def\BNE_scanop_a #1#2%
{%
    \expandafter\BNE_scanop_b\expandafter#1\romannumeral`&&@#2%
}%
\def\BNE_scanop_b #1#2%
{%
    \unless\ifcat#2\relax
           \ifcsname BNE_itself_#1#2\endcsname
           \BNE_scanop_c
    \fi\fi
    \BNE_foundop_a #1#2%
}%
\def\BNE_scanop_c #1#2#3#4#5% #1#2=\fi\fi
{%
    #1#2%
    \expandafter\BNE_scanop_d\csname BNE_itself_#4#5\expandafter\endcsname
    \romannumeral`&&@%
}%
\def\BNE_scanop_d #1#2%
{%
    \unless\ifcat#2\relax
           \ifcsname BNE_itself_#1#2\endcsname
           \BNE_scanop_c
    \fi\fi
    \BNE_foundop #1#2%
}%
%    \end{macrocode}
% If a postfix say |?s| is defined and |?r| is encountered the |?| will have
% been interpreted as a shortcut to |?s| and then the |r| will be found with
% the parser (after having executed the already found postfix) now looking for
% another operator so the error message will be |Operator? (got `r')| which is
% doubly confusing... well, let's not dwell on that.
%    \begin{macrocode}
\def\BNE_foundop_a #1%
{%
    \ifcsname BNE_precedence_#1\endcsname
        \csname BNE_precedence_#1\expandafter\endcsname
        \expandafter #1%
    \else
        \expandafter\BNE_getop_a\romannumeral`&&@%
        \xint_afterfi{\XINT_expandableerror
        {Operator? (got `#1'). Hit I<RET><operator>}}%
     \fi
}%
\def\BNE_foundop #1{\csname BNE_precedence_#1\endcsname #1}%
%    \end{macrocode}
% \subsection{Expansion spanning; opening and closing parentheses}\leavevmode
%    \begin{macrocode}
\def\BNE_tmpa #1#2#3#4#5%
{%
    \def#1% start
    {%
        \expandafter#2\romannumeral`&&@\BNE_getnext
    }%
    \def#2##1% check
    {%
        \xint_UDsignfork
          ##1{\expandafter#3\romannumeral`&&@#4}%
            -{#3##1}%
        \krof
    }%
    \def#3##1##2% checkp
    {%
        \ifcase ##1%
           \expandafter\BNE_done
        \or\expandafter#5%
        \else
           \expandafter#3\romannumeral`&&@\csname BNE_op_##2\expandafter\endcsname
        \fi
    }%
    \def#5%
    {%
        \XINT_expandableerror
        {An extra ) has been removed. Hit <RET>, fingers crossed.}%
        \expandafter#2\romannumeral`&&@\expandafter\BNE_put_op_first
        \romannumeral`&&@\BNE_getop_legacy
    }%
}%
\let\BNE_done\space
\def\BNE_getop_legacy #1%
{%
    \expanded{\unexpanded{{#1}}\expandafter}\romannumeral`&&@\BNE_getop
}%
\expandafter\BNE_tmpa
    \csname BNE_start\expandafter\endcsname
    \csname BNE_check\expandafter\endcsname
    \csname BNE_checkp\expandafter\endcsname
    \csname BNE_op_-xii\expandafter\endcsname
    \csname BNE_extra_)\endcsname
\catcode`) 11
\def\BNE_tmpa #1#2#3#4#5#6%
{%
    \def #1##1% op_(
    {%
        \expandafter #4\romannumeral`&&@\BNE_getnext
    }%
    \def #2##1% op_)
    {%
        \expanded{\unexpanded{\BNE_put_op_first{##1}}\expandafter}\romannumeral`&&@\BNE_getop
    }%
    \def #3% oparen
    {%
        \expandafter #4\romannumeral`&&@\BNE_getnext
    }%
    \def #4##1% check-
    {%
        \xint_UDsignfork
            ##1{\expandafter#5\romannumeral`&&@#6}%
              -{#5##1}%
        \krof
    }%
    \def #5##1##2% checkp
    {%
        \ifcase ##1\expandafter\BNE_missing_)
        \or \csname BNE_op_##2\expandafter\endcsname
        \else
          \expandafter #5\romannumeral`&&@\csname BNE_op_##2\expandafter\endcsname
        \fi
    }%
}%
\expandafter\BNE_tmpa
    \csname BNE_op_(\expandafter\endcsname
    \csname BNE_op_)\expandafter\endcsname
    \csname BNE_oparen\expandafter\endcsname
    \csname BNE_check-_)\expandafter\endcsname
    \csname BNE_checkp_)\expandafter\endcsname
    \csname BNE_op_-xii\endcsname
\let\BNE_precedence_)\xint_c_i
\def\BNE_missing_) 
   {\XINT_expandableerror{Missing ). Hit <RET> to proceed}%
    \xint_c_ \BNE_done }%
\catcode`) 12
%    \end{macrocode}
% \subsection{The comma as binary operator}
% At |1.4|, it is simply a union operator for 1D oples. Inserting directly
% here a |<comma><space>| separator (as in earlier releases) in accumulated
% result would avoid having to do it on output but to the cost of diverging
% from \xintexprname upstream code, and to have to let the |\evaltohex| output
% routine handle comma separated values rather than braced values.
%    \begin{macrocode}
\def\BNE_tmpa #1#2#3#4#5%
{%
    \def #1##1% \BNE_op_,
    {%
      \expanded{\unexpanded{#2{##1}}\expandafter}%
      \romannumeral`&&@\expandafter#3\romannumeral`&&@\BNE_getnext
    }%
    \def #2##1##2##3##4{##2##3{##1##4}}% \BNE_exec_,
    \def #3##1% \BNE_check-_,
    {%
      \xint_UDsignfork
        ##1{\expandafter#4\romannumeral`&&@#5}%
          -{#4##1}%
      \krof
    }%
    \def #4##1##2% \BNE_checkp_,
    {%
      \ifnum ##1>\xint_c_iii
        \expandafter#4%
           \romannumeral`&&@\csname BNE_op_##2\expandafter\endcsname
      \else
        \expandafter##1\expandafter##2%
      \fi
    }%
}%
\expandafter\BNE_tmpa
    \csname BNE_op_,\expandafter\endcsname
    \csname BNE_exec_,\expandafter\endcsname
    \csname BNE_check-_,\expandafter\endcsname
    \csname BNE_checkp_,\expandafter\endcsname
    \csname BNE_op_-xii\endcsname
\expandafter\let\csname BNE_precedence_,\endcsname\xint_c_iii
%    \end{macrocode}
% \subsection{The minus as prefix operator of variable precedence level}
% This |\BNE_Op_opp| caused trouble at |1.4| as it must be \emph{f}-expandable,
% whereas earlier it expanded inside |\csname...\endcsname| context, so I
% could define it as \centeredline{|\if-#1\else\if0#10\else-#1\fi\fi|} where
% |#1| was the
% first token of unbraced argument but this meant at |1.4| an added
% |\xint_firstofone| here. Well let's return to sanity at |1.4a| and not add
% the |\xint_firstofone| and simply default |\BNE_Op_opp| to |\xintiiOpp|,
% which it should have been all along! And on this occasion let's trim user
% documentation of complications.
%
% The package used to need to define unary minus operator with precedences 12,
% 14, and 18.  It also defined it at level 16 but this was unneedeed actually,
% no operator possibly generating usage of an |op_-xvi|.
%
% At |1.5| the right precedence of powers was lowered to 17, so we now need
% here only 12, 14, and 17.
%
% Due to |\bnumdefinfix| it is needed to support also, perhaps, the other
% levels 13, 15, 16, 18, ....  This will be done only if necessary and is the
% reason why the macros |\BNE_defminus_a| and |\BNE_defminus_b| are given
% permanent names. In fact it is now |\BNE_defbin_b| which will decide to invoke
% or not the |\BNE_defminus_a|, and we activate it here only for the base
% precedence 12.
%
% The |\XINT_global| are inexistent in upstream at |1.4f| as it does not
% incorporate yet some analog to |\bnumdefinfix/\bnumdefpostfix|.
%    \begin{macrocode}
\def\BNE_defminus_b #1#2#3#4#5%
{%
    \XINT_global\def #1% \BNE_op_-<level>
    {%
      \expandafter #2\romannumeral`&&@\expandafter#3%
      \romannumeral`&&@\BNE_getnext
    }%
    \XINT_global\def #2##1##2##3% \BNE_exec_-<level>
    {%
      \expandafter ##1\expandafter ##2\expandafter
       {\expandafter{\romannumeral`&&@\BNE_Op_opp##3}}%
    }%
    \XINT_global\def #3##1% \BNE_check-_-<level>
    {%
      \xint_UDsignfork
        ##1{\expandafter #4\romannumeral`&&@#1}%
          -{#4##1}%
      \krof
    }%
    \XINT_global\def #4##1##2% \BNE_checkp_-<level>
    {%
      \ifnum ##1>#5%
        \expandafter #4%
        \romannumeral`&&@\csname BNE_op_##2\expandafter\endcsname
      \else
        \expandafter ##1\expandafter ##2%
      \fi
    }%
}%
\def\BNE_defminus_a #1%
{%
    \expandafter\BNE_defminus_b
    \csname BNE_op_-#1\expandafter\endcsname
    \csname BNE_exec_-#1\expandafter\endcsname
    \csname BNE_check-_-#1\expandafter\endcsname
    \csname BNE_checkp_-#1\expandafter\endcsname
    \csname xint_c_#1\endcsname
}%
\BNE_defminus_a {xii}%
%    \end{macrocode}
% \subsection{The infix operators.}
% I could have at the |1.4| refactoring injected usage of |\expanded| here,
% but kept in sync with upstream \xintexprname code. Any \emph{x}-expandable
% macro can easily be converted into an \emph{f}-expandable one using
% |\expanded|, so this is no serious limitation.
%
% Macro names are somewhat bad and there is much risk of confusion in future
% maintenance of |\BNE_Op_| prefix (used for |\BNE_Op_add| etc...; besides
% this should have been |\BNE_Op_Add|) and |\BNE_op_| prefix (used for
% |\BNE_op_+| etc...).
%
% At |1.5| decision is made to anticipate the announced upstream change to let
% the power operators be right associative, matching Python behaviour. This
% change is simply implemented by hardcoding in |\BNE_checkp_<op>| the right
% precedence which so far, for such operators, had been identical with the
% left precedence (upstream has examples of direct coding without
% formalization). In fact the right precedence existed already as argument to
% |\BNE_defbin_b| as the precedence to assign to unary minus following |<op>|.
%
% Note1: although it is easy to change the left precedence at user level, the
% right precedence is now more inaccessible. But on the other hand \bnumname
% provides |\bnumdefinfix| so all is customizable at user level.
%
% Note2: Tacit multiplication is not really a separate operator, it is
% the |*| with an elevated left precedence, which costs nothing to create and
% this precedence is stored in chardef token |\BNE_prec_tacit|.
%
% Compared to upstream, we use here numbers as arguments to |\BNE_defbin_b|,
% and convert to roman numerals internally, also the operator macro is passed
% as a control sequence not as its name (and |#6| and |#7| are permuted in
% |\BNE_defbin_c|).
%    \begin{macrocode}
\def\BNE_defbin_c #1#2#3#4#5#6#7%
{%
  \XINT_global\def #1##1% \BNE_op_<op>
  {%
    \expanded{\unexpanded{#2{##1}}\expandafter}%
    \romannumeral`&&@\expandafter#3\romannumeral`&&@\BNE_getnext
  }%
  \XINT_global\def #2##1##2##3##4% \BNE_exec_<op>
  {%
    \expandafter##2\expandafter##3\expandafter
      {\expandafter{\romannumeral`&&@#7##1##4}}%
  }%
  \XINT_global\def #3##1% \BNE_check-_<op>
  {%
    \xint_UDsignfork
      ##1{\expandafter#4\romannumeral`&&@#5}%
        -{#4##1}%
    \krof
  }%
  \XINT_global\def #4##1##2% \BNE_checkp_<op>
  {%
    \ifnum ##1>#6%
      \expandafter#4%
      \romannumeral`&&@\csname BNE_op_##2\expandafter\endcsname
    \else 
      \expandafter ##1\expandafter ##2%
    \fi
  }%
}%
\def\BNE_defbin_b #1#2#3#4%
{%
    \expandafter\BNE_defbin_c
    \csname BNE_op_#1\expandafter\endcsname
    \csname BNE_exec_#1\expandafter\endcsname
    \csname BNE_check-_#1\expandafter\endcsname
    \csname BNE_checkp_#1\expandafter\endcsname
    \csname BNE_op_-\romannumeral\ifnum#3>12 #3\else 12\fi
            \expandafter\endcsname
    \csname xint_c_\romannumeral#3\endcsname #4%
  \XINT_global
    \expandafter
    \let\csname BNE_precedence_#1\expandafter\endcsname
        \csname xint_c_\romannumeral#2\endcsname
    \unless
    \ifcsname BNE_exec_-\romannumeral\ifnum#3>12 #3\else 12\fi\endcsname
%    \end{macrocode}
% This will execute only for |#3>12| as |\BNE_exec_-xii| exists.
%    \begin{macrocode}
     \expandafter\BNE_defminus_a\expandafter{\romannumeral#3}%
    \fi
}%
\BNE_defbin_b  +   {12} {12}  \BNE_Op_add
\BNE_defbin_b  -   {12} {12}  \BNE_Op_sub
\BNE_defbin_b  *   {14} {14}  \BNE_Op_mul
\BNE_defbin_b  /   {14} {14}  \BNE_Op_divround
\BNE_defbin_b {//} {14} {14}  \BNE_Op_div
\BNE_defbin_b {/:} {14} {14}  \BNE_Op_mod
\BNE_defbin_b  ^   {18} {17}  \BNE_Op_pow
%    \end{macrocode}
% At upstream, we can use shortcut
% \centeredline{|\expandafter\def\csname BNE_itself_**\endcsname {^}|}
% but it means then that any redefinition of |^| propagates to |**|, besides it
% creates a special case which would need consideration by
% |\BNE_dotheitselves|, or special restrictions to add to user documentation.
% Better to simply handle |**| as a full operator.
%    \begin{macrocode}
\BNE_defbin_b {**} {18} {17}  \BNE_Op_pow
\expandafter\def\csname BNE_itself_**\endcsname {**}%
\expandafter\def\csname BNE_itself_//\endcsname {//}%
\expandafter\def\csname BNE_itself_/:\endcsname {/:}%
\let\BNE_prec_tacit\xint_c_xvi
%    \end{macrocode}
% \subsection{\csh{bnumdefinfix}: extending the syntax}\leavevmode
% |#1| gives the operator characters, |#2| the associated macro, |#3| its
% left-precedence and |#4| its right precedence (as integers).
%
% The "itself" definitions are done in such a way that unambiguous abbreviations
% work; but in case of ambiguity the first defined operator is used.
%
% However, if for example operator |$a| was defined after |$ab|, then although
% |$| will use |$ab| which was defined first, |$a| will use as expected the
% second defined operator.
%
% The mismatch |\BNE_defminus_a| vs |\BNE_defbin_b| is inherited from
% upstream, I keep it to simplify maintenance.
%    \begin{macrocode}
\def\bnumdefinfix #1#2#3#4%
{%
    \edef\BNE_tmpa{#1}%
    \edef\BNE_tmpa{\xint_zapspaces_o\BNE_tmpa}%
    \edef\BNE_tmpL{\the\numexpr#3\relax}%
    \edef\BNE_tmpL{\ifnum\BNE_tmpL<4 4\else\ifnum\BNE_tmpL<23 \BNE_tmpL\else 22\fi\fi}%
    \edef\BNE_tmpR{\the\numexpr#4\relax}%
    \edef\BNE_tmpR{\ifnum\BNE_tmpR<4 4\else\ifnum\BNE_tmpR<23 \BNE_tmpR\else 22\fi\fi}%
    \BNE_defbin_b \BNE_tmpa\BNE_tmpL\BNE_tmpR #2%
    \expandafter\BNE_dotheitselves\BNE_tmpa\relax
  \ifxintverbose
    \PackageInfo{bnumexpr}{infix operator \BNE_tmpa\space
    \ifxintglobaldefs globally \fi
        does
        \unexpanded{#2}\MessageBreak with precedences \BNE_tmpL, \BNE_tmpR;}%
  \fi
}%
\def\BNE_dotheitselves#1#2%
{%
    \if#2\relax\expandafter\xint_gobble_ii
    \else
  \XINT_global
      \expandafter\edef\csname BNE_itself_#1#2\endcsname{#1#2}%
      \unless\ifcsname BNE_precedence_#1\endcsname
  \XINT_global
        \expandafter\edef\csname BNE_precedence_#1\endcsname
                         {\csname BNE_precedence_\BNE_tmpa\endcsname}%
  \XINT_global
        \expandafter\odef\csname BNE_op_#1\endcsname
                         {\csname BNE_op_\BNE_tmpa\endcsname}%
      \fi
    \fi
    \BNE_dotheitselves{#1#2}%
}%
%    \end{macrocode}
% \subsection{\csh{bnumdefpostfix}}
% Support macros for postfix operators only need to be \emph{x}-expandable.
%    \begin{macrocode}
\def\bnumdefpostfix #1#2#3%
{%
    \edef\BNE_tmpa{#1}%
    \edef\BNE_tmpa{\xint_zapspaces_o\BNE_tmpa}%
    \edef\BNE_tmpL{\the\numexpr#3\relax}%
    \edef\BNE_tmpL{\ifnum\BNE_tmpL<4 4\else\ifnum\BNE_tmpL<23 \BNE_tmpL\else 22\fi\fi}%
  \XINT_global
    \expandafter\let\csname BNE_precedence_\BNE_tmpa\expandafter\endcsname
                    \csname xint_c_\romannumeral\BNE_tmpL\endcsname
  \XINT_global
    \expandafter\def\csname BNE_op_\BNE_tmpa\endcsname ##1%
    {%
        \expandafter\BNE_put_op_first
        \expanded{{{#2##1}}\expandafter}\romannumeral`&&@\BNE_getop
    }%
    \expandafter\BNE_dotheitselves\BNE_tmpa\relax
  \ifxintverbose
    \PackageInfo{bnumexpr}{postfix operator \BNE_tmpa\space
    \ifxintglobaldefs globally \fi
        does \unexpanded{#2}\MessageBreak
        with precedence \BNE_tmpL;}%
  \fi
}%
%    \end{macrocode}
% \subsection{! as postfix factorial operator}\leavevmode
%    \begin{macrocode}
\bnumdefpostfix{!}{\BNE_Op_fac}{20}%
%    \end{macrocode}
% \subsection{Cleanup}\leavevmode
%    \begin{macrocode}
\let\BNEtmpa\relax
\let\BNE_tmpa\relax \let\BNE_tmpb\relax \let\BNE_tmpc\relax
\let\BNE_tmpR\relax \let\BNE_tmpL\relax
\BNErestorecatcodes%
%    \end{macrocode}
% \MakePercentComment
%</package>
%<*dtx>
\CheckSum {1205}%
\makeatletter\check@checksum\makeatother%
\Finale%
%% End of file xint.dtx
